\documentclass[11pt]{article}
\usepackage[margin=0.8in]{geometry}
\usepackage{booktabs,longtable,array}
\usepackage{graphicx}
\usepackage{float}
\usepackage{pdflscape}
\usepackage{siunitx}
\usepackage[T1]{fontenc}
\usepackage{lmodern}
\usepackage{hyperref}
\sisetup{detect-all}
\renewcommand{\thetable}{S\arabic{table}}
\renewcommand{\thefigure}{S\arabic{figure}}
\setlength{\LTleft}{0pt}
\setlength{\LTright}{0pt}

\title{Supporting Information for\\
BioInteract: Interpretable Drug--Target Interaction Prediction via Residue-Level Cross-Attention with Protein-Language and Physicochemical Priors}
\author{ }
\date{ }

\begin{document}
\maketitle

\noindent This document reports values recoverable from the released Davis
records, configuration files, saved result files, source code, and the fixed
BindingDB external-validation manifest included with this revision. ``Classifier
score'' denotes the sigmoid output for the binary high-affinity class; it is not
an experimental affinity. Model-native attention and Grad-CAM values are
attribution scores, not contact assignments.

\section*{Model configuration and input features}

\begin{table}[H]
\centering
\caption{Model architecture confirmed from the checkpoint-embedded model configurations.}
\begin{tabular}{ll}
\toprule
Component & Archived setting \\
\midrule
Drug encoder & GINE; 3 layers; hidden dimension 256; dropout 0.20 \\
Atom and bond input & 52 node features; 16 edge features \\
Protein language model & \texttt{esm2\_t30\_150M\_UR50D}; 640-dimensional residues \\
Protein projection & 256-dimensional projection; 32-dimensional domain-label embedding \\
Cross-attention & 8 heads; attention dimension 256; dropout 0.10 \\
Pooling & Gated \\
Prediction head & Hidden dimensions 256 and 128; dropout 0.30; classification task \\
\bottomrule
\end{tabular}
\end{table}

\begin{table}[H]
\centering
\caption{Protein-residue input channels and their availability.}
\begin{tabular}{p{0.26\linewidth}p{0.66\linewidth}}
\toprule
Channel & Archived implementation and limitation \\
\midrule
ESM-2 & Frozen 640-dimensional residue embeddings; proteins are truncated at 1,200 residues. \\
Physicochemical vector & Four values per residue: hydrophobicity, normalized molecular weight, charge, and polarity. \\
Domain-label channel & The configuration reserves a 32-dimensional embedding. The label builder uses \texttt{NONE} when no per-protein JSON annotation is available. \\
Annotation availability & No \texttt{data/domain\_annotations/*.json} files are present in the released package. Thus the reproducible Davis input uses the unknown label at every residue and is not curated Pfam/InterPro information. \\
\bottomrule
\end{tabular}
\end{table}

\begin{table}[H]
\centering
\caption{Training provenance for the historical entity-split checkpoints and the
separately trained exact-sequence-grouped strict models.}
\small
\begin{tabular}{p{0.31\linewidth}p{0.61\linewidth}}
\toprule
Setting & Value \\
\midrule
Historical entity-split scope & Model architecture only; the three historical entity-split checkpoints do not retain optimizer, scheduler, warmup, or loss state \\
Historical Random and Drug-ID logs & Final-run headers record weighted BCE, label smoothing 0.05, AdamW ($10^{-4}$, $10^{-5}$), and cosine warmup \\
Historical Target-ID trace & Older \texttt{run\_split.py} header; only loss and validation AUROC are logged, so optimizer, warmup, and smoothing are not recoverable from the artifact \\
Historical threshold protocol & A validation F1 threshold is frozen before each canonical test evaluation \\
Strict-model initialisation and data & Each exact-sequence-grouped model was newly initialised with seed 42 and trained only on its new Davis training partition, using the same preprocessing and architecture \\
Strict-model optimisation & Weighted BCE with 0.05 label smoothing; AdamW (learning rate $10^{-4}$, weight decay $10^{-5}$); batch size 64; gradient-norm clipping 1.0; five-epoch linear warm-up followed by cosine decay; maximum 100 epochs \\
Strict-model selection & Highest validation AUROC with patience 15; validation F1 threshold selected after loading that checkpoint and frozen for test evaluation \\
Strict cold-target artifact & \texttt{revision/analysis/sequence\_grouped\_cold\_target.pt}; best epoch 24 of 39; validation AUROC 0.911568; threshold 0.837581; SHA-256 {\scriptsize\texttt{8ea58e1f056b5e33c309787880a811c797393510c459d9442b5dbe67e632fe5c}} \\
Strict cold-both artifact & \texttt{revision/analysis/sequence\_grouped\_cold\_both.pt}; best epoch 6 of 21; validation AUROC 0.910630; threshold 0.935704; SHA-256 {\scriptsize\texttt{8b7c0f374983a35163a7d3243c89732504a5a97ef4d10c557b2841a627873972}} \\
\bottomrule
\end{tabular}
\end{table}

\begin{table}[H]
\centering
\caption{Computational environment documented during final reproducibility audit.}
\begin{tabular}{p{0.30\linewidth}p{0.62\linewidth}}
\toprule
Item & Value \\
\midrule
Python & 3.10.14 \\
PyTorch & 2.4.0+cu121 \\
PyTorch Geometric & 2.6.1 \\
RDKit & 2025.3.6 \\
\texttt{fair-esm} & 2.0.0 \\
Biopython and RapidFuzz & 1.87 and 3.14.5 \\
CUDA availability during audit & Available \\
Audit GPU & NVIDIA GeForce RTX 4080 \\
Strict-split run record & \texttt{device: cuda} in \texttt{strict\_split\_metrics.json} \\
Historical checkpoint files & \texttt{best.pt}, \texttt{best\_random.pt}, \texttt{best\_cold\_target.pt}, \texttt{best\_cold\_drug.pt} \\
Strict-model checkpoint files & \texttt{revision/analysis/sequence\_grouped\_cold\_target.pt} and \texttt{revision/analysis/sequence\_grouped\_cold\_both.pt} (hashes in Table~S3) \\
\bottomrule
\end{tabular}
\end{table}

\clearpage

\section*{Dataset, partitions, and metrics}

\begin{table}[H]
\centering
\caption{Canonical current-release Davis test-partition summaries. Values are
from the full-precision checkpoint reconstruction, frozen validation thresholds,
and the associated prediction files.}
\begin{tabular}{lrrrrrrr}
\toprule
Partition & Test pairs & Positives & AUROC & AUPRC & F1 & Precision & Recall \\
\midrule
Random & 6,011 & 276 & 0.904 & 0.560 & 0.565 & 0.602 & 0.533 \\
Target-ID-held-out & 5,984 & 250 & 0.930 & 0.525 & 0.534 & 0.517 & 0.552 \\
Drug-ID-held-out & 5,746 & 245 & 0.733 & 0.167 & 0.100 & 0.273 & 0.061 \\
\bottomrule
\end{tabular}
\end{table}

\begin{table}[H]
\centering
\caption{Bootstrap intervals for the canonical current-release entity-based partitions.
Each interval uses 600 ordinary pair-level bootstrap replicates sampled with
replacement at the test-pair level, NumPy seed 42, and the 2.5th and 97.5th
percentiles of valid replicate metrics. Replicates containing only one class
are discarded. The main figure shows point estimates only.}
\begin{tabular}{lcc}
\toprule
Partition & AUROC 95\% interval & AUPRC 95\% interval \\
\midrule
Random & [0.881, 0.924] & [0.498, 0.624] \\
Target-ID-held-out & [0.911, 0.948] & [0.454, 0.584] \\
Drug-ID-held-out & [0.703, 0.763] & [0.129, 0.211] \\
\bottomrule
\end{tabular}
\end{table}

\begin{table}[H]
\centering
\caption{Exact-sequence-grouped strict within-Davis evaluations from
\texttt{strict\_split\_metrics.json}. The two models were trained from
initialisation under the protocol in Table~S3. Threshold-dependent quantities use
the corresponding validation-set threshold; all strict-split entries are point
estimates, not bootstrap confidence intervals.}
\scriptsize
\textbf{Pair accounting}\par\smallskip
\begin{tabular}{lrrrrr}
\toprule
Protocol & Train pairs & Validation pairs & Test pairs & Test positives & Excluded cross-block pairs \\
\midrule
Exact-sequence-grouped cold-target & 21,080 & 2,992 & 5,984 & 208 & 0 \\
Exact-sequence-grouped cold-both & 15,190 & 264 & 1,144 & 38 & 13,458 \\
\bottomrule
\end{tabular}

\medskip
\textbf{Test metrics}\par\smallskip
\begin{tabular}{lrrrrrr}
\toprule
Protocol & AUROC & AUPRC & Threshold & F1 & Precision & Recall \\
\midrule
Exact-sequence-grouped cold-target & 0.873 & 0.383 & 0.838 & 0.468 & 0.517 & 0.428 \\
Exact-sequence-grouped cold-both & 0.552 & 0.038 & 0.936 & 0.000 & 0.000 & 0.000 \\
\bottomrule
\end{tabular}

\medskip
\raggedright For the cold-both protocol, drug identifiers and exact
target-sequence groups are partitioned independently at 70/10/20. Only pairs in
the corresponding train--train, validation--validation, and test--test blocks
are retained; the remaining cross-block pairs are excluded. The cold-both
validation block contains 10 positive pairs.
\end{table}

\section*{Reproducibility details for revision analyses}

\noindent\textbf{Identifier-partition similarity audits.}
Drug similarities use RDKit~2025.3.6 Morgan fingerprints (radius 2, 1,024 bits,
\texttt{includeChirality=False}) and Tanimoto similarity. Protein similarities
for the original Target-ID-held-out audit use RapidFuzz~3.14.5
\texttt{fuzz.ratio}/100. For sequences $x,y$, this is the normalised Indel
similarity $1-d_{\mathrm{Indel}}/(|x|+|y|)$, where substitutions are represented
as one deletion plus one insertion.

\medskip
\noindent\textbf{Strict-test global alignment.}
Biopython~1.87 \texttt{Align.PairwiseAligner} is used in global mode with match
$=1$, mismatch $=0$, gap-opening $=-1$, and gap-extension $=-0.5$. Full-length
identity is the number of exact-residue matches divided by alignment length; the
maximum across all training targets is retained for each held-out target.

\medskip
\noindent\textbf{Bootstrap intervals.}
The saved 95\% intervals use ordinary, unstratified pair-level resampling with
replacement, 600 requested replicates, and NumPy random seed 42. Each replicate
contains the same number of test pairs as the corresponding test set; one-class
replicates are skipped. Bounds are the 2.5th and 97.5th percentiles of the valid
AUROC or AUPRC replicate distribution.

\medskip
\noindent\textbf{Grad-CAM.}
Hooks are attached to \texttt{drug\_encoder.layers.2}, the third and final GINE
message-passing layer. Gradients are taken from the pre-sigmoid positive-class
logit, pooled across graph nodes to obtain channel weights, rectified, and then
normalised separately within each molecular graph.

\medskip
\noindent\textbf{Attention sparsity.}
For pair $q$, residue scores $s_{qj}=\sum_i M_{qij}$ are normalised as
$\widetilde{s}_{qj}=s_{qj}/\max_k s_{qk}$. The reported fraction is
$|\{(q,j):\widetilde{s}_{qj}<0.1\}|/\sum_q L_q$. It uses all 1,506 positive Davis
pairs with the archived \texttt{checkpoints/best.pt}, not a single test set.
Every pair is normalised before the 1,301,380 valid residue scores are
concatenated; thus original training, validation, and test memberships are all
included and longer targets contribute more residue-level observations.
The archived flat-score array is
\texttt{BioInteract/results/figure\_data/attention\_distribution.npz}; it is
the sole source of the attention-distribution figure and cumulative curve in the
main text.

\begin{center}
\small
\textbf{Residue-attribution summary.}\par\smallskip
\begin{tabular}{lr}
\toprule
Quantity & Value \\
\midrule
Positive Davis pairs & 1,506 \\
Residue-score observations & 1,301,380 \\
Mean / standard deviation & 0.004351 / 0.046243 \\
Median & 0.000272 \\
95th / 99th percentile & 0.003617 / 0.057193 \\
Fraction below 0.1 & 0.992319 \\
Attribution mass in top 1\% / 5\% scores & 0.788528 / 0.904465 \\
\bottomrule
\end{tabular}
\end{center}

\begin{center}
\small
\textbf{Exact-sequence-grouped target diagnostic controls.} All values use the 5,984-pair strict target test set with 208 positives.\\[3pt]
\begin{tabular}{lrr}
\toprule
Control & AUROC & AUPRC \\
\midrule
BioInteract & 0.873 & 0.383 \\
Drug-only Morgan logistic regression & 0.764 & 0.164 \\
Drug-promiscuity prior & 0.764 & 0.163 \\
Target-only constant prior & 0.500 & 0.035 \\
Label-shuffled drug-only control & 0.407 & 0.029 \\
\bottomrule
\end{tabular}
\end{center}

\clearpage

\section*{Feature and input audits}

\begin{table}[H]
\centering
\caption{Atom and bond feature specification from \texttt{src/data/mol\_graph.py}.}
\begin{tabular}{llr}
\toprule
Level & Feature family & Dimensions \\
\midrule
Atom & Atom type one-hot (20 choices including unknown) & 20 \\
Atom & Hybridization, formal charge, and hydrogen-count one-hot vectors & 18 \\
Atom & Degree, aromaticity, and ring membership & 3 \\
Atom & Ring-size membership for sizes 3--8 & 6 \\
Atom & Donor/acceptor, Crippen logP and MR, Gasteiger charge & 5 \\
Atom total &  & 52 \\
Bond & Bond type, conjugation, ring membership, stereo, and direction & 16 \\
\bottomrule
\end{tabular}
\end{table}

\begin{longtable}{@{}l@{\quad}l@{\quad}r@{\quad}c@{\quad}r@{\quad}r@{}}
\caption{Davis nominal ABL1 input audit. Every listed record has the same complete SHA-256 digest {\scriptsize\texttt{cd6a78b82fede215cdb80a797f228098a41178fdde0753728f3781e6221fb1ee}}; values in the last columns are the archived D0010 and D0017 affinities in nM. The nominal labels do not alter the model sequence input.}\\
\toprule
Target ID & Nominal name & Length & Matches ABL1 & D0010 & D0017 \\
\midrule
\endfirsthead
\toprule
Target ID & Nominal name & Length & Matches ABL1 & D0010 & D0017 \\
\midrule
\endhead
T0001 & ABL1(E255K) & 1167 & Yes & 0.047 & 0.047 \\
T0002 & ABL1(F317I) & 1167 & Yes & 0.63 & 0.10 \\
T0003 & ABL1(F317I)p & 1167 & Yes & 0.18 & 0.041 \\
T0004 & ABL1(F317L) & 1167 & Yes & 0.11 & 0.032 \\
T0005 & ABL1(F317L)p & 1167 & Yes & 0.029 & 0.019 \\
T0006 & ABL1(H396P) & 1167 & Yes & 0.062 & 0.025 \\
T0007 & ABL1(H396P)p & 1167 & Yes & 0.057 & 0.046 \\
T0008 & ABL1(M351T) & 1167 & Yes & 0.037 & 0.016 \\
T0009 & ABL1(Q252H) & 1167 & Yes & 0.086 & 0.037 \\
T0010 & ABL1(Q252H)p & 1167 & Yes & 0.039 & 0.064 \\
T0011 & ABL1(T315I) & 1167 & Yes & 21 & 890 \\
T0012 & ABL1(T315I)p & 1167 & Yes & 3.6 & 120 \\
T0013 & ABL1(Y253F) & 1167 & Yes & 0.036 & 0.058 \\
T0014 & ABL1 & 1167 & Yes & 0.12 & 0.029 \\
T0015 & ABL1p & 1167 & Yes & 0.057 & 0.046 \\
\bottomrule
\end{longtable}

\begin{table}[H]
\centering
\caption{Saved non-mutant EGFR--D0014 functional-group attribution values.}
\begin{tabular}{lr}
\toprule
Functional group & Attribution \\
\midrule
Carbonyl & 0.431 \\
Amide & 0.389 \\
Ether & 0.136 \\
Halogen & 0.134 \\
Amino & 0.090 \\
Aromatic ring & 0.000 \\
Heterocycle N & 0.000 \\
\bottomrule
\end{tabular}
\end{table}

\begin{table}[H]
\centering
\caption{Saved non-mutant BRAF--D0023 failure-mode summary. The archived affinity is positive at the study threshold, whereas the classifier score is low. The score and attributions are model outputs, not an explanation of binding.}
\begin{tabular}{lr}
\toprule
Quantity & Value \\
\midrule
High-affinity-class score & 0.099 \\
Archived Davis affinity (nM) & 0.19 \\
Aromatic-ring attribution & 0.059 \\
Hydroxyl attribution & 0.000 \\
Heterocycle-N attribution & 0.000 \\
\bottomrule
\end{tabular}
\end{table}

\begin{table}[H]
\centering
\caption{Trainable parameter count recovered from \texttt{interpretability\_report.json}.}
\begin{tabular}{lr}
\toprule
Quantity & Value \\
\midrule
Total trainable parameters & 2,442,083 \\
\bottomrule
\end{tabular}
\end{table}

\clearpage
\section*{Strict external BindingDB evaluation}

% Source identifier: BindingDB_BindingDB_Articles_202607
Table~S13 records the frozen external evaluation requested by the reviewers.
The official curated-articles snapshot was handled independently of Davis, then
filtered under the pre-specified exact-$K_d$ and exact double-novel protocol.
This is an external exact-entity test, not a claim of scaffold or remote-homology
novelty, and no BindingDB label was used to retrain, calibrate, select the
checkpoint, or select the threshold.

\begin{table}[H]
\centering
\small
\caption{Strict BindingDB external-validation provenance, curation, and frozen
result. The source archive was \texttt{BindingDB\_BindingDB\_Articles\_202607\_tsv.zip}
(SHA-256 \texttt{d2584d1519318d00ab5f46289da5ab3549affe732d598a5072f8777b6b3b5262}).
Only finite exact $K_d$ measurements in nM for valid single-chain proteins
(maximum 1,200 residues) and valid ligand structures were considered. Exact
canonical-isomeric-SMILES and full-normalised-sequence matches to every Davis
ligand or target were excluded. Eligible measurements were grouped by canonical
ligand--sequence pair; threshold-conflicting groups were removed at the pair
level, and every remaining pair retained its median exact $K_d$. Confidence
intervals are pair-stratified bootstrap 95\% intervals from 600 replicates
(seed 42).}
\label{tab:bindingdb_external_si}
\begin{tabular}{p{0.42\linewidth}p{0.49\linewidth}}
\toprule
Item & Value \\
\midrule
Source snapshot & {\scriptsize\texttt{BindingDB\_BindingDB\_\linebreak[0]Articles\_202607}} (official BindingDB curated articles TSV) \\
Source rows & 93,712 \\
Final positives / negatives & 516 / 1,416 \\
Unique ligands / targets & 1,195 / 396 \\
Sequential row-level exclusions & 82,988 censored/non-exact $K_d$; 8,261 multi-chain; 139 invalid sequences; 8 invalid SMILES; 64 Davis-ligand matches; 94 Davis-sequence matches \\
Eligible measurements before pair reconciliation & 2,158 \\
Threshold-conflicting replicate groups & 15 ligand--target pairs, comprising 44 measurement rows; conflicts are counted at the pair level \\
Threshold-consistent measurement rows & 2,114 \\
Median replicate aggregation & 143 retained pairs had $\geq2$ consistent measurements; 182 repeated measurement rows were collapsed by the within-pair median $K_d$ \\
Final external pairs & 1,932 \\
Frozen model and threshold & \texttt{checkpoints/best\_random.pt}; 0.5959881544 (existing Davis random-validation threshold) \\
Protein features & Newly extracted \texttt{fair-esm} 2.0.0 \texttt{esm2\_t30\_150M\_UR50D}; all 396 caches passed 640-dimensional and length checks \\
AUROC (95\% CI) & 0.5597 [0.5313, 0.5885] \\
AUPRC (95\% CI) & 0.3404 [0.3138, 0.3724] \\
Frozen-threshold F1 / precision / recall & 0.0701 / 0.7308 / 0.0368 \\
Two-run reproducibility & Identical per-pair probabilities; the complete SHA-256 is recorded in \texttt{bindingdb\_external\_manifest.json} \\
Interpretation & Limited external transfer; the result does not support a broad DTI-transfer claim and is not used for structural or attribution interpretation. \\
\bottomrule
\end{tabular}
\end{table}

\clearpage
\section*{Breadth check across model families}

The main-text comparison in Table~1 is a single-variable ladder: the same
molecular graph encoder throughout, first without a protein language model,
then with one, then with cross-attention in place of concatenation. Table~S14
places that ladder alongside two further architectures trained by the same
harness, under the same protocol described in Table~S16, on the random split.
The two additional models are a sequence convolutional model in the DeepDTA
family and the low-rank bilinear attention and bilinear pooling of DrugBAN
applied to the same two encoders used by BioInteract. They were run on the
random split only, as a breadth check across families rather than as a second
controlled contrast, and their protocol-dependence was therefore not assessed.

The five architectures fall within roughly 0.03 AUROC of one another on this
split: the sequence convolutional model, which uses neither a molecular graph
nor a protein language model, reaches AUROC 0.912 and AUPRC 0.569, the bilinear
interaction module reaches 0.907 and 0.547, and BioInteract reaches 0.913 and
0.545. This is the basis for the statement in the main text that the random
split separates architectures weakly, so that a favourable random-split number
is poor evidence for a design choice and the informative contrasts are the
entity-held-out ones.

\begin{table}[H]
\centering
\small
\caption{Random-split breadth check. All values were measured in this study
under the shared protocol of Table~S16. The first three rows are the main-text
ladder; the last two are the additional architectures. ``PLM'' indicates
whether the model consumes the frozen \texttt{esm2\_t30\_150M\_UR50D} residue
representation. Best value in each column is in bold.}
\begin{tabular}{llccc}
\toprule
Model & Interaction paradigm & PLM & AUROC & AUPRC \\
\midrule
GraphDTA & Graph + sequence CNN & No & 0.885 & 0.358 \\
ESM2-GraphConcat & Graph + PLM, concatenation & Yes & 0.907 & 0.478 \\
\textbf{BioInteract (ours)} & Graph + PLM, cross-attention & Yes & \textbf{0.913} & 0.545 \\
DeepDTA & Sequence CNN & No & 0.912 & \textbf{0.569} \\
ESM2-Bilinear & Graph + PLM, bilinear attention & Yes & 0.907 & 0.547 \\
\bottomrule
\end{tabular}
\label{tab:breadth_si}
\end{table}

\clearpage
\section*{Identity-stratified cold-target results}

Table~S15 gives the numerical values plotted in main-text Figure~3. Held-out
targets in the exact-sequence-grouped cold-target partition are grouped by their
maximum full-length global sequence identity to any training target, computed
with the Biopython global aligner (match $=1$, mismatch $=0$, gap opening $=-1$,
gap extension $=-0.5$); identity is exact-residue matches divided by full
alignment length. The two outer strata contain 75 and 3 positive pairs
respectively, so their AUPRC estimates are unstable; the ordering across strata
is not monotone and should not be read as a homology-transfer trend.

\begin{table}[H]
\centering
\small
\caption{BioInteract performance on the exact-sequence-grouped cold-target test,
stratified by maximum global sequence identity to the training targets. The
positive prevalence of the full cold-target test is 0.035.}
\begin{tabular}{lrrrrrr}
\toprule
Maximum identity & Test targets & Test pairs & Positive pairs &
Positive prevalence & AUROC & AUPRC \\
\midrule
$<0.40$         & 31 & 2{,}108 & 75 & 0.036 & 0.806 & 0.305 \\
$0.40$--$<0.60$ & 28 & 1{,}904 & 77 & 0.040 & 0.908 & 0.456 \\
$0.60$--$<0.80$ & 24 & 1{,}632 & 53 & 0.032 & 0.908 & 0.466 \\
$\geq 0.80$     &  5 &    340 &  3 & 0.009 & 0.862 & 0.121 \\
\midrule
All             & 88 & 5{,}984 & 208 & 0.035 & 0.873 & 0.383 \\
\bottomrule
\end{tabular}
\label{tab:identity_bins_si}
\end{table}

\clearpage
\section*{Matched-protocol comparison provenance}

Table~S16 records one row per training run behind main-text Table~1. Every run
used the same Davis records, the same partition for its protocol, seed 42,
weighted binary cross-entropy with 0.05 label smoothing, AdamW (learning rate
$10^{-4}$, weight decay $10^{-5}$), an effective batch size of 64 assembled from
micro-batches of 16 with gradient accumulation, automatic mixed precision,
gradient-norm clipping at 1.0, a five-epoch linear warm-up followed by cosine
decay, a maximum of 30 epochs, and early stopping at patience 6 on validation
AUROC. Validation and test inference were run in full precision. The reported
threshold was selected once on that run's validation predictions by maximising
F1 and then frozen for the test evaluation. Each configuration was trained once;
repeated-seed variability was not estimated, so small differences between
adjacent rows should not be interpreted.

\begin{table}[H]
\centering
\scriptsize
\caption{Per-run provenance for the matched-protocol comparison. ``Best epoch''
is the epoch whose validation AUROC was highest and whose weights were restored
before the test evaluation.}
\resizebox{\textwidth}{!}{%
\begin{tabular}{llrrrrrrr}
\toprule
Model & Protocol & Parameters & Epochs & Best epoch & Val.\ AUROC &
Threshold & Test AUROC & Test AUPRC \\
\midrule
BioInteract & random & 2{,}442{,}083 & 21 & 15 & 0.9456 & 0.9036 & 0.9128 & 0.5454 \\
DeepDTA & random & 919{,}169 & 30 & 27 & 0.9316 & 0.9207 & 0.9125 & 0.5692 \\
ESM2-Bilinear & random & 1{,}533{,}697 & 30 & 28 & 0.9175 & 0.8270 & 0.9065 & 0.5470 \\
ESM2-GraphConcat & random & 1{,}335{,}553 & 30 & 29 & 0.9236 & 0.8277 & 0.9071 & 0.4785 \\
GraphDTA & random & 2{,}264{,}513 & 26 & 20 & 0.8895 & 0.8822 & 0.8851 & 0.3581 \\
BioInteract & target\_id\_held\_out & 2{,}442{,}083 & 25 & 19 & 0.9214 & 0.8867 & 0.9340 & 0.5693 \\
DeepDTA & target\_id\_held\_out & 919{,}169 & 30 & 26 & 0.8846 & 0.7149 & 0.9136 & 0.4932 \\
ESM2-GraphConcat & target\_id\_held\_out & 1{,}335{,}553 & 23 & 17 & 0.8553 & 0.8958 & 0.8780 & 0.3687 \\
GraphDTA & target\_id\_held\_out & 2{,}264{,}513 & 22 & 16 & 0.8457 & 0.6843 & 0.8557 & 0.3056 \\
BioInteract & drug\_id\_held\_out & 2{,}442{,}083 & 7 & 1 & 0.5708 & 0.6939 & 0.7126 & 0.1259 \\
ESM2-GraphConcat & drug\_id\_held\_out & 1{,}335{,}553 & 7 & 1 & 0.5833 & 0.5038 & 0.6836 & 0.0968 \\
GraphDTA & drug\_id\_held\_out & 2{,}264{,}513 & 7 & 1 & 0.6180 & 0.6050 & 0.6793 & 0.1050 \\
\bottomrule
\end{tabular}
}
\label{tab:matched_provenance_si}
\end{table}

\clearpage
\section*{BindingDB filtering waterfall and snapshot provenance}

Table~S17 records the provenance of the two BindingDB snapshots audited in this
revision and the reconciliation of the two filter orderings. Both orderings
apply exactly the same criteria and both retain the same number of eligible
measurements for a given snapshot; they differ only in the order in which the
measurement-quality and target-validity tests are applied, which changes how
individual rows are attributed between the first three counters. The
measurement-first ordering used in main-text Table~4 was adopted because it
isolates the dominant cause of the reduction, namely rows that report no $K_d$
at all.

\begin{table}[H]
\centering
\small
\caption{Provenance of the two BindingDB snapshots and reconciliation of the two
audit orderings. The curated-articles snapshot is the pre-specified source of
the frozen external evaluation; the complete release was added in this revision
as an independent verification that the filter definitions scale.}
\begin{tabular}{p{0.30\linewidth}p{0.31\linewidth}p{0.31\linewidth}}
\toprule
Item & Curated articles & Complete release \\
\midrule
Source version & {\scriptsize\texttt{BindingDB\_BindingDB\_\linebreak[0]Articles\_202607}} & \texttt{BindingDB\_All\_202608} \\
Archive file & {\scriptsize\texttt{BindingDB\_BindingDB\_\linebreak[0]Articles\_202607\_tsv.zip}} & \texttt{\scriptsize BindingDB\_All\_202608\_tsv.zip} \\
Archive bytes & 18{,}114{,}757 & 592{,}986{,}963 \\
Archive SHA-256 & {\scriptsize\texttt{d2584d1519318d00ab5f46289da5ab35\linebreak[0]49affe732d598a5072f8777b6b3b5262}} & {\scriptsize\texttt{8dc30541d668e5e403ba1fe9a682b76c\linebreak[0]42a50e1711aac44329c62ba24ccd57db}} \\
TSV member & {\scriptsize\texttt{BindingDB\_BindingDB\_\linebreak[0]Articles.tsv}} & \texttt{BindingDB\_All.tsv} \\
Input rows & 93{,}712 & 3{,}234{,}499 \\
Rows with no $K_d$ reported & 91{,}062 (97.2\%) & 3{,}104{,}331 (96.0\%) \\
Eligible measurements, measurement-first ordering & 2{,}158 & 52{,}085 \\
Eligible measurements, single-chain-first ordering & 2{,}158 & 52{,}085 \\
Pair-level threshold conflicts removed & 15 pairs / 44 rows & 1{,}992 pairs / 5{,}233 rows \\
Retained double-novel pairs & 1{,}932 & 38{,}193 \\
Retained positives / negatives & 516 / 1{,}416 & 13{,}264 / 24{,}929 \\
Unique ligands / targets & 1{,}195 / 396 & 26{,}433 / 3{,}462 \\
Frozen AUROC (95\% CI) & 0.560 [0.531, 0.589] & 0.533 [0.526, 0.539] \\
Frozen AUPRC (95\% CI) & 0.340 [0.314, 0.372] & 0.383 [0.378, 0.390] \\
Frozen-threshold F1 / precision / recall & 0.070 / 0.731 / 0.037 & 0.044 / 0.513 / 0.023 \\
Prediction SHA-256 & {\scriptsize\texttt{9e50dfadbdcb9974d6bc2dcefbdc4edb\linebreak[0]f4bd089b14598d76724da48fa40bba1c}} & {\scriptsize\texttt{303ba15ab0e605607102d08e7b2fc3c4\linebreak[0]cee888d11cc531d5aced121d00cf9013}} \\
Repeat-run check & Two runs gave identical per-pair probabilities & Evaluated once; the repeat run was not performed \\
Role in this study & Pre-specified frozen external evaluation & Verification of filter scaling and a larger frozen external evaluation \\
\bottomrule
\end{tabular}
\label{tab:bindingdb_waterfall_si}
\end{table}

\end{document}
